\documentclass[12pt]{article}

\usepackage{url,verbatim, amsmath, amssymb, amsfonts, dsfont, color}
\textwidth 15true cm
\textheight 9true in
\oddsidemargin 0.30true in
\evensidemargin 0.30true in
\topmargin -0.3true in
\headsep 0.4true in

\def\sgn{{\rm sign}}
\def\drightarrow{{\stackrel{d}{\rightarrow}}}
\def\dotsim{\stackrel{.}{\sim}}
\def\goesd{\stackrel{d}{\rightarrow}}
\def\goesp{\stackrel{p}{\rightarrow}}
\def\goeslaw{\stackrel{{\cal L}}{\longrightarrow}}
%
\def\vp{{\varphi}}
\def\E{{\rm E}}
\def\Pvalue{{$p$-value}}
\def\Pvalues{{$p$-values}}
\def\pr{{\rm pr}}
\def\Pr{{\rm pr}}
\def\logit{{\rm logit}}
\def\T{{ \mathrm{\scriptscriptstyle T} }}
\def\diag{{\rm diag}}  
\def\half{{\textstyle{1\over 2}}}
\def\ith{{$i{\rm th}$ }}
\def\var{{\rm var}}
\def\cl{{{\rm c}\ell}}
\def\checkx{{\checked\hspace{-6pt}\setminus}}
\def\grad{{\nabla}}
\newcommand{\Var}{\text{Var}}

\sloppy %permits line-breaking of urls

\def\blue{\color{blue}}
\def\red{\color{red}}
\def\sgn{{\rm sign}}
\def\drightarrow{{\stackrel{d}{\rightarrow}}}
\begin{document}

\noindent{\bf Questions Week 3} {\hfill STA 2212S 2026}\\

%{\bf Due January 14} \hfill{\it MS, Exercise 5.2}

\bigskip

\begin{enumerate}

\item For $\theta\in\mathbb{R}$, show that if $g(X;\theta)$ is an unbiased estimating function under the model for $X$, i.e.$$\int g(x;\theta)f(x;\theta)dx = 0,$$
that $$\frac{\var\{g(X;\theta)\}}{\E\{-\partial g(X;\theta)/\partial\theta\}^2} \ge I_1^{-1}(\theta),$$
where $I_1(\theta) = \var\{\partial \log f(X;\theta)/\partial\theta \}$ is the expected Fisher information in one observation. This is a generalization of the Cram\'{e}r-Rao lower bound.

If this question is too easy, you can show for $\theta\in\mathbb{R}^p$, that 
$$
\E\left\{\frac{\partial g(X;\theta)}{\partial\theta^T}\right\}^{-1} \var\{g(X;\theta)\} \E\left\{\frac{\partial g(X;\theta)^T}{\partial\theta}\right\}^{-1} - I^{-1}(\theta)
$$ is non-negative definite. Here $g(X;\theta)$ is a $p\times 1$ vector, still unbiased for $0$, and we assume $\E\left\{\frac{\partial g(X;\theta)}{\partial\theta^T}\right\}$ is invertible. But I won't ask you to prove this part on the midterm.

\textcolor{red}{Differentiate 
\begin{align*}
    &\frac{\partial }{\partial \theta}\int g(x; \theta) f(x ; \theta) dx = 0\\
    &\int \frac{\partial}{\partial \theta}g(x ; \theta) f(x ; \theta) dx = 0\\
    &\int \frac{\partial g(x ; \theta)}{\partial \theta}f(x ; \theta) + g(x ; \theta) \frac{\partial f(x ; \theta)}{\partial \theta} dx = 0
\end{align*}
Now, 
\begin{equation*}
    E\left[\frac{\partial g(X ; \theta)}{\partial \theta}\right] = \int \frac{\partial g(x ; \theta)}{\partial \theta}f(x ; \theta) dx 
\end{equation*}
Also, notice that 
\begin{equation*}
    \frac{\partial f(x ; \theta)}{\partial \theta} = f(x;\theta)\frac{\partial \ell(\theta ; x)}{\partial x}
\end{equation*}
where $\ell(\theta ; x) = \log f(x; \theta)$ is the log-likelihood. Thus,
\begin{equation*}
    \int g(x;\theta)\frac{\partial f(x;\theta)}{\partial \theta}dx = \int g(x;\theta)\frac{\partial \ell(\theta;x)}{\partial \theta} f(x;\theta) dx = E\left[g(X;\theta) \frac{\partial \ell(\theta;X)}{\partial \theta}\right]
\end{equation*}
hence
\begin{equation*}
    E\left[g(X;\theta) \frac{\partial \ell(\theta;X)}{\partial \theta}\right] = -E\left[\frac{\partial g(X ; \theta)}{\partial \theta}\right]
\end{equation*}
By the Cauchy-Schwarz inequality, 
\begin{equation*}
    \left(E\left[g(X;\theta) \frac{\partial \ell(\theta;X)}{\partial \theta}\right]\right)^2 \leq E(g(X;\theta)^2)I_1(\theta) 
\end{equation*}
thus 
\begin{equation*}
    E(g(X;\theta^2))I_1(\theta) \geq E\left[-\frac{\partial g(X;\theta)}{\partial \theta}\right]^2
\end{equation*}
Since $g(X;\theta)$ is unbiased, then $E(g(X;\theta)^2) = \Var(g(X;\theta))$, thus 
\begin{equation*}
    \frac{\Var(g(X;\theta))}{E\left[-\frac{\partial g(X;\theta)}{\partial \theta}\right]^2} \leq I_1^{-1}(\theta)
\end{equation*}
follows immediately. 
}
\item ({\it MS Problem 7.1}) Suppose that $X_1, \dots, X_n$ are i.i.d.~normally distributed random variables with unknown mean $\mu$ and variance $\sigma^2$. 
Using the pivot $(n-1)S^2/\sigma^2$, where $$S^2 = \frac{1}{n-1}\sum_{i=1}^n(X_i-\bar X)^2,$$ we can obtain a 95\% confidence interval $[k_1S^2,k_2S^2]$ for some constants $k_1$ and $k_2$. Find expressions for $k_1$ and $k_2$ if this confidence interval has minimum length. Evaluate $k_1$ and $k_2$ when $n=10$.

\textcolor{red}{
Since the $X_i$ are i.i.d. normal, then 
\begin{equation*}
    \frac{(n-1)S^2}{\sigma^2} \sim \chi_{n-1}^2
\end{equation*}
For a 95\% confidence interval $[a, b]$, we want to have $a, b$ such that 
\begin{equation*}
    P\left( a \leq \frac{(n-1)S^2 }{\sigma^2} \leq b \right) = 0.95
\end{equation*}
Rewriting yields 
\begin{equation*}
    P\left(\frac{(n-1)S^2}{b} \leq \sigma^2 \leq \frac{(n-1)S^2}{a}\right) = 0.95
\end{equation*}
thus we take $k_1 = (n-1)S^2/b$ and $k_2 = (n-1)S^2 / a$ and optimize the distance $k_2 - k_1$. Notice that 
\begin{equation*}
    k_2 - k_1 = (n-1)S^2 \left(\frac1a - \frac 1b\right) \propto \frac 1a - \frac 1b
\end{equation*}
Let $f(a, b) = \frac 1a - \frac 1b$. We wish to minimize $f(a, b)$ subject to $F(b) - F(a) = 0.95$ where $F$ is the cdf of $\chi_{n-1}^2$. By Lagrangian multipliers, 
\begin{equation*}
    \frac{1}{a^2} = \lambda f(a)  \qquad \frac{1}{b^2} = \lambda f(b)
\end{equation*}
which implies 
\begin{equation*}
    a^2 f(a) = b^2f(b)
\end{equation*}
thus the minimal 95\% CI is obtained from solving a system $a^2f(a) = b^2f(b)$ and $F(b) - F(a) = 0.95$ and picking $k_1 = (n-1)S^2/b$ and $k_2 = (n-1)S^2/a$. \\
When $n = 10$, then $k_1 = 0.345$ and $k_2 = 2.741$. 
}
\end{enumerate}
\end{document}

